% SEGMENT OLP-0129-S01
% Part: first-order-logic
% Chapter: completeness
% Section: complete-sets

% Definition of complete consistent sets. Properties of complete sets required
% for completeness proved are in provability.tex in the
% chapter on the proof system used.

\documentclass[../../../include/open-logic-section]{subfiles}

\begin{document}

\iftag{FOL}
      {\olfileid{fol}{com}{ccs}}
      {\olfileid{pl}{com}{ccs}}

\olsection{நிறைவான முரணற்ற \usetoken{P}{sentence}-கணங்கள்}

% SEGMENT OLP-0129-S02
\begin{defn}[நிறைவான கணம்]
\ollabel{def:complete-set} !!{sentence}s கொண்ட கணம்~$\Gamma$,
ஒவ்வொரு !!{sentence}~$!A$ க்கும் $!A \in \Gamma$ அல்லது
$\lnot !A \in \Gamma$ எனில், அப்போதும் அப்போது மட்டுமே
\emph{!!{complete}} ஆனது.
\end{defn}

% SEGMENT OLP-0129-S03
\begin{explain}
!!^{complete} வாக்கியக் கணங்கள் எந்தக் கேள்வியையும் விடையளிக்காமல் விடுவதில்லை.
ஒவ்வொரு !!{sentence}~$!A$ க்கும், $!A$ மெய்யா பொய்யா என்பதை~$\Gamma$
``கூறுகிறது.'' நிறைவு தேற்றத்தின் நிறுவலுக்கு அப்பாலும் !!{complete} கணங்கள்
முக்கியமானவை. எடுத்துக்காட்டாக, !!{complete} ஆகவும் அடிகோளாக்கத்தக்கதாகவும்
உள்ள ஒரு கோட்பாடு எப்போதும் தீர்மானிக்கத்தக்கது.
\end{explain}

% SEGMENT OLP-0129-S04
\begin{explain}
!!^{complete} முரணற்ற கணங்கள் நிறைவு நிறுவலில் முக்கியமானவை. ஏனெனில்,
முரணற்ற !!{sentence}s கொண்ட ஒவ்வொரு கணம்~$\Gamma$ உம் !!a{complete} முரணற்ற
கணம்~$\Gamma^*$ இல் அடங்கியுள்ளது என்பதை உறுதிசெய்ய முடியும்.
!!^a{complete} முரணற்ற கணம் ஒவ்வொரு !!{sentence}~$!A$ க்கும் $!A$ அல்லது
அதன் மறுப்பு~$\lnot !A$ ஐக் கொண்டிருக்கும்; இரண்டையும் கொண்டிருக்காது.
குறிப்பாக, \iftag{FOL}{அணு !!{sentence}s}{கூற்று மாறிகள்} ஒவ்வொன்றுக்கும்
இது பொருந்தும். ஆகவே, \iftag{FOL}{!!{constant}s ஐ ஏற்றவாறு சேர்த்து
விரிவாக்கிய மொழியிலுள்ள}{} !!a{complete} முரணற்ற கணத்திலிருந்து,
\iftag{FOL}{!!{predicate}s இன் பொருள்விளக்கம்}{கூற்று மாறிகளுக்கு ஒதுக்கப்படும்
மெய்மதிப்பு} $\Gamma^*$ விலுள்ள \iftag{FOL}{அணு !!{sentence}s}{!!{propositional variable}s}
எவை என்பதன்படி வரையறுக்கப்படும் ஒரு
\iftag{FOL}{!!a{structure}}{!!a{valuation}} ஐக் கட்டமைக்கலாம். இந்த
\iftag{FOL}{!!{structure}}{!!{valuation}} ஆனது~$\Gamma^*$ விலுள்ள எல்லா
!!{sentence}s ஐயும் (ஆகவே~$\Gamma$ விலுள்ள அனைத்தையும்) மெய்யாக்கும் என்று
பின்னர் காட்டலாம். இதை நிறுவ, $\lnot !A \in \Gamma^*$ அப்போதும் அப்போது
மட்டுமே $!A \notin \Gamma^*$, $(!A \lor !B) \in \Gamma^*$ அப்போதும்
அப்போது மட்டுமே $!A \in \Gamma^*$ அல்லது $!B \in \Gamma^*$, இதுபோன்ற
பண்புகள் தேவைப்படும்.
\end{explain}

% SEGMENT OLP-0129-S05
அடுத்து வருவதில் $\Proves$ இன் தற்சுட்டுப் பண்பு, ஒருபோக்குத்தன்மை,
கடப்புப் பண்பு ஆகியவற்றை மறைமுகமாகப் பலமுறை பயன்படுத்துவோம் (காண்க
\iftag{FOL}{%
  \tagrefs{prfSC/{fol:seq:ptn:sec},prfND/{fol:ntd:ptn:sec},prfAX/{fol:axd:ptn:sec},prfTab/{fol:tab:ptn:sec}}}{%
  \tagrefs{prfSC/{pl:seq:ptn:sec},prfND/{pl:ntd:ptn:sec},prfAX/{pl:axd:ptn:sec},prfTab/{pl:tab:ptn:sec}}}).

% SEGMENT OLP-0129-S06
\begin{prop}
\ollabel{prop:ccs}
$\Gamma$ ஆனது !!{complete} ஆகவும் முரணற்றதாகவும் உள்ளது என்க. அப்போது:
\begin{enumerate}
\item \ollabel{prop:ccs-prov-in} $\Gamma \Proves !A$ எனில்,
  $!A \in \Gamma$.

\tagitem{prvAnd}{\ollabel{prop:ccs-and} $!A \land !B \in \Gamma$
  அப்போதும் அப்போது மட்டுமே $!A \in \Gamma$, $!B \in \Gamma$ இரண்டும்.}{}

\tagitem{prvOr}{\ollabel{prop:ccs-or} $!A \lor !B \in \Gamma$ அப்போதும்
  அப்போது மட்டுமே $!A \in \Gamma$ அல்லது $!B \in \Gamma$.}{}

\tagitem{prvIf}{\ollabel{prop:ccs-if} $!A \lif !B \in \Gamma$ அப்போதும்
  அப்போது மட்டுமே $!A \notin \Gamma$ அல்லது $!B \in \Gamma$.}{}
\end{enumerate}
\end{prop}

% SEGMENT OLP-0129-S07
\begin{proof}
பின்வரும் அனைத்திலும் $\Gamma$ ஆனது !!{complete} ஆகவும் முரணற்றதாகவும் உள்ளது
எனக் கொள்வோம்.
\begin{enumerate}
\item $\Gamma \Proves !A$ எனில், $!A \in \Gamma$.

$\Gamma \Proves !A$ என்க. இதற்கு மாறாக $!A \notin \Gamma$ எனக் கொள்வோம்.
$\Gamma$ ஆனது !!{complete} ஆதலால், $\lnot !A \in \Gamma$. பின்வரும்
மேற்கோளின்படி,
\iftag{FOL}{%
  \tagrefs{prfSC/{fol:seq:prv:prop:explicit-inc},
    prfND/{fol:ntd:prv:prop:explicit-inc},
    prfAX/{fol:axd:prv:prop:explicit-inc},
    prfTab/{fol:tab:prv:prop:explicit-inc}}}{%
  \tagrefs{prfSC/{pl:seq:prv:prop:explicit-inc},
    prfND/{pl:ntd:prv:prop:explicit-inc},
    prfAX/{pl:axd:prv:prop:explicit-inc},
    prfTab/{pl:tab:prv:prop:explicit-inc}}},
$\Gamma$ முரணுடையது. இது $\Gamma$ முரணற்றது என்ற கருதுகோளுக்கு முரணாகும்.
எனவே $!A \notin \Gamma$ என்பது இயலாது; ஆகவே $!A \in \Gamma$.

% SEGMENT OLP-0129-S08
\tagitem{defAnd}{}{%
\iftag{probAnd}{பயிற்சி.}{%
$!A \land !B \in \Gamma$ அப்போதும் அப்போது மட்டுமே $!A \in \Gamma$,
$!B \in \Gamma$ இரண்டும்:

முன்னோக்கிய திசைக்கு $!A \land !B \in \Gamma$ என்க. அப்போது
\iftag{FOL}{%
  \tagrefs{prfSC/{fol:seq:ppr:prop:provability-land},
    prfND/{fol:ntd:ppr:prop:provability-land},
    prfAX/{fol:axd:ppr:prop:provability-land},
    prfTab/{fol:tab:ppr:prop:provability-land}}}{%
  \tagrefs{prfSC/{pl:seq:ppr:prop:provability-land},
    prfND/{pl:ntd:ppr:prop:provability-land},
    prfAX/{pl:axd:ppr:prop:provability-land},
    prfTab/{pl:tab:ppr:prop:provability-land}}}, உருப்படி~(1) படி,
$\Gamma \Proves !A$, $\Gamma \Proves !B$. \olref{prop:ccs-prov-in} படி,
தேவையானவாறு $!A \in \Gamma$, $!B \in \Gamma$.

மறுதிசைக்கு $!A \in \Gamma$, $!B \in \Gamma$ என்க. பின்வரும்
மேற்கோளின் உருப்படி~(2) படி,
\iftag{FOL}{%
  \tagrefs{prfSC/{fol:seq:ppr:prop:provability-land},
    prfND/{fol:ntd:ppr:prop:provability-land},
    prfAX/{fol:axd:ppr:prop:provability-land},
    prfTab/{fol:tab:ppr:prop:provability-land}}}{%
  \tagrefs{prfSC/{pl:seq:ppr:prop:provability-land},
    prfND/{pl:ntd:ppr:prop:provability-land},
    prfAX/{pl:axd:ppr:prop:provability-land},
    prfTab/{pl:tab:ppr:prop:provability-land}}},
$\Gamma \Proves !A \land !B$. \olref{prop:ccs-prov-in} படி,
$!A \land !B \in \Gamma$.}}

% SEGMENT OLP-0129-S09
\tagitem{defOr}{}{%
\iftag{probOr}{பயிற்சி.}{%
முதலில், $!A \lor !B \in \Gamma$ எனில் $!A \in \Gamma$ அல்லது
$!B \in \Gamma$ என்பதைக் காட்டுவோம். $!A \lor !B \in \Gamma$ என்றும்,
ஆனால் $!A \notin \Gamma$, $!B \notin \Gamma$ என்றும் கொள்வோம்.
$\Gamma$ ஆனது !!{complete} ஆதலால், $\lnot !A \in \Gamma$,
$\lnot !B \in \Gamma$. பின்வரும் மேற்கோளின் உருப்படி~(1) படி,
\iftag{FOL}{%
  \tagrefs{prfSC/{fol:seq:ppr:prop:provability-lor},
    prfND/{fol:ntd:ppr:prop:provability-lor},
    prfAX/{fol:axd:ppr:prop:provability-lor},
    prfTab/{fol:tab:ppr:prop:provability-lor}}}{%
  \tagrefs{prfSC/{pl:seq:ppr:prop:provability-lor},
    prfND/{pl:ntd:ppr:prop:provability-lor},
    prfAX/{pl:axd:ppr:prop:provability-lor},
    prfTab/{pl:tab:ppr:prop:provability-lor}}},
$\Gamma$ முரணுடையது; இது ஒரு முரண்பாடு. எனவே $!A \in \Gamma$ அல்லது
$!B \in \Gamma$.

மறுதிசைக்கு $!A \in \Gamma$ அல்லது $!B \in \Gamma$ என்க. பின்வரும்
மேற்கோளின் உருப்படி~(2) படி,
\iftag{FOL}{%
  \tagrefs{prfSC/{fol:seq:ppr:prop:provability-lor},
    prfND/{fol:ntd:ppr:prop:provability-lor},
    prfAX/{fol:axd:ppr:prop:provability-lor},
    prfTab/{fol:tab:ppr:prop:provability-lor}}}{%
  \tagrefs{prfSC/{pl:seq:ppr:prop:provability-lor},
    prfND/{pl:ntd:ppr:prop:provability-lor},
    prfAX/{pl:axd:ppr:prop:provability-lor},
    prfTab/{pl:tab:ppr:prop:provability-lor}}},
$\Gamma \Proves !A \lor !B$. \olref{prop:ccs-prov-in} படி,
தேவையானவாறு $!A \lor !B \in \Gamma$.}}

% SEGMENT OLP-0129-S10
\tagitem{defIf}{}{%
\iftag{probIf}{பயிற்சி.}{%
முன்னோக்கிய திசைக்கு $!A \lif !B \in \Gamma$ என்க. இதற்கு மாறாக
$!A \in \Gamma$, $!B \notin \Gamma$ எனக் கொள்வோம். இந்தக்
கருதுகோள்களின்படி $!A \lif !B \in \Gamma$, $!A \in \Gamma$.
பின்வரும் மேற்கோளின் உருப்படி~(1) படி,
\iftag{FOL}{%
  \tagrefs{prfSC/{fol:seq:ppr:prop:provability-lif},
    prfND/{fol:ntd:ppr:prop:provability-lif},
    prfAX/{fol:axd:ppr:prop:provability-lif},
    prfTab/{fol:tab:ppr:prop:provability-lif}}}{%
  \tagrefs{prfSC/{pl:seq:ppr:prop:provability-lif},
    prfND/{pl:ntd:ppr:prop:provability-lif},
    prfAX/{pl:axd:ppr:prop:provability-lif},
    prfTab/{pl:tab:ppr:prop:provability-lif}}},
$\Gamma \Proves !B$. ஆனால் \olref{prop:ccs-prov-in} படி
$!B \in \Gamma$; இது $!B \notin \Gamma$ என்ற கருதுகோளுக்கு முரணாகும்.

மறுதிசையில் முதலில் $!A \notin \Gamma$ என்ற நிலையை எடுத்துக்கொள்வோம்.
$\Gamma$ ஆனது !!{complete} ஆதலால் $\lnot !A \in \Gamma$. பின்வரும்
மேற்கோளின் உருப்படி~(2) படி,
\iftag{FOL}{%
  \tagrefs{prfSC/{fol:seq:ppr:prop:provability-lif},
    prfND/{fol:ntd:ppr:prop:provability-lif},
    prfAX/{fol:axd:ppr:prop:provability-lif},
    prfTab/{fol:tab:ppr:prop:provability-lif}}}{%
  \tagrefs{prfSC/{pl:seq:ppr:prop:provability-lif},
    prfND/{pl:ntd:ppr:prop:provability-lif},
    prfAX/{pl:axd:ppr:prop:provability-lif},
    prfTab/{pl:tab:ppr:prop:provability-lif}}},
$\Gamma \Proves !A \lif !B$. மீண்டும் \olref{prop:ccs-prov-in} படி,
தேவையானவாறு $!A \lif !B \in \Gamma$.

இப்போது $!B \in \Gamma$ என்ற நிலையை எடுத்துக்கொள்வோம். பின்வரும்
மேற்கோளின் உருப்படி~(2) ஐ மீண்டும் பயன்படுத்தினால்,
\iftag{FOL}{%
  \tagrefs{prfSC/{fol:seq:ppr:prop:provability-lif},
    prfND/{fol:ntd:ppr:prop:provability-lif},
    prfTab/{fol:tab:ppr:prop:provability-lif},
    prfAX/{fol:axd:ppr:prop:provability-lif}}}{%
  \tagrefs{prfSC/{pl:seq:ppr:prop:provability-lif},
    prfND/{pl:ntd:ppr:prop:provability-lif},
    prfTab/{pl:tab:ppr:prop:provability-lif},
    prfAX/{pl:axd:ppr:prop:provability-lif}}},
$\Gamma \Proves !A \lif !B$. \olref{prop:ccs-prov-in} படி,
$!A \lif !B \in \Gamma$.}}
\end{enumerate}
\end{proof}

% SEGMENT OLP-0129-S11
\tagprob{FOL}
\begin{prob}
\olref[fol][com][ccs]{prop:ccs} இன் நிறுவலை முடிக்கவும்.
\end{prob}
\tagendprob

\tagprob{notFOL}
\begin{prob}
\olref[pl][com][ccs]{prop:ccs} இன் நிறுவலை முடிக்கவும்.
\end{prob}
\tagendprob

\end{document}
